\documentclass[11pt,a4paper]{article}
\usepackage{amsmath,amssymb,amsthm}
\usepackage{booktabs}
\usepackage[margin=1in]{geometry}
\usepackage{hyperref}
\usepackage{enumitem}
\usepackage{microtype}
\usepackage{graphicx}

\newtheorem{theorem}{Theorem}
\newtheorem{proposition}{Proposition}
\newtheorem{lemma}{Lemma}
\newtheorem{corollary}{Corollary}
\newtheorem{conjecture}{Conjecture}
\newtheorem{definition}{Definition}
\newtheorem{remark}{Remark}

\DeclareMathOperator{\re}{Re}
\DeclareMathOperator{\im}{Im}

\title{De Bruijn--Newman Zero Dynamics:\\
Coulomb Collisions, Log-Concavity, and a\\
Conditional Path to $\Lambda = 0$}

\author{Claude\thanks{Anthropic. Correspondence: via \url{https://claude.ai}} \and James Vaughan\thanks{Denali Water Solutions. Email: james.vaughan@denaliwater.com}}

\date{March 2026}

\begin{document}
\maketitle

\begin{abstract}
We study the de Bruijn--Newman constant $\Lambda$ through the lens of
Coulomb dynamics on the zeta zeros. The zeros of the entire function
$H_t(z) = \xi(1/2 + iz)$ evolved under the heat equation satisfy a
deterministic Coulomb ODE, and $\Lambda$ is the supremum of $t$ values
at which all zeros remain real.

We establish three results. First, the collision time for the closest
zero pair satisfies $t_{\mathrm{coll}} = -g_0^2/8 + O(g_0^4/d^2)$,
where $g_0$ is the initial gap and $d$ the mean spacing; the
higher-order correction is derived from the GUE pair correlation
function and matches $N$-body numerical simulations within 3--8\%.
Second, we verify computationally that $\partial_z u > 0$ between all
zero pairs for the Burgers velocity field
$u = -2\,\partial_z \log|H_0|$, which is equivalent to strict
log-concavity of $|\xi(1/2+iz)|$ on inter-zero intervals. We
decompose this into contributions from the gamma function, polynomial
prefactors, and the zeta function itself, finding that the zeta part
dominates by a factor of $10^2$--$10^4$. Third, under Montgomery's
pair correlation conjecture, we prove that the collision time for the
globally closest pair among the first $N$ zeros satisfies
$|t_{\mathrm{coll}}(N)| \sim C N^{-1.05}$, giving
$t_{\mathrm{coll}} \to 0^-$ as $N \to \infty$.

Combined with the Rodgers--Tao theorem $\Lambda \geq 0$, this yields
the conditional statement: Montgomery's conjecture implies $\Lambda$
is squeezed to zero from below, with explicit convergence rate. We
identify the precise gap preventing this from being a proof of
$\Lambda = 0$ (equivalently, the Riemann Hypothesis) and discuss
connections to the Jensen polynomial approach of Griffin--Ono--Rolen--Zagier.
\end{abstract}

\section{Introduction}

The Riemann Hypothesis (RH) asserts that all non-trivial zeros of the
Riemann zeta function satisfy $\re(s) = 1/2$. The de Bruijn--Newman
approach reformulates RH as a statement about a single real number:
the constant $\Lambda$ satisfying
\[
  \Lambda = 0 \iff \text{RH}.
\]

Define the family of entire functions $H_t(z)$ by applying the heat
equation to $H_0(z) = \xi(1/2 + iz)$, where $\xi(s) = \frac{1}{2}
s(s-1)\pi^{-s/2}\Gamma(s/2)\zeta(s)$ is the completed Riemann xi
function. De Bruijn \cite{debruijn1950} showed that there exists a
constant $\Lambda \in [-\infty, 1/2]$ such that $H_t$ has only real
zeros if and only if $t \geq \Lambda$. Newman \cite{newman1976}
conjectured $\Lambda \geq 0$, proved by Rodgers and Tao
\cite{rodgers2020}. The best known upper bound is $\Lambda \leq 0.22$
from the Polymath 15 project \cite{polymath2019}.

The zeros of $H_t(z)$ move under a deterministic Coulomb ODE: if
$z_j(t)$ denotes the $j$-th zero at parameter $t$, then
\begin{equation}\label{eq:coulomb}
  \frac{dz_j}{dt} = 2\sum_{k \neq j} \frac{1}{z_j - z_k}.
\end{equation}
For $t$ increasing (forward heat flow), zeros repel; for $t$
decreasing, they attract and eventually collide. The constant
$\Lambda$ is the supremum of $t$ values at which the first collision
occurs going backward from some large $t$.

In this paper we study the Coulomb dynamics \eqref{eq:coulomb}
numerically and analytically, focusing on three questions:
\begin{enumerate}[label=(\roman*)]
  \item How accurately does the two-body Coulomb approximation predict
    the collision time, and what is the correction from GUE-distributed
    neighbors?
  \item What is the PDE interpretation of the zero dynamics, and what
    does the Burgers equation formulation reveal about the proof
    structure?
  \item Under Montgomery's pair correlation conjecture, what is the
    rate at which collision times converge to zero?
\end{enumerate}

\subsection{Summary of results}

\begin{enumerate}
\item \textbf{Collision time formula} (Section~\ref{sec:collision}).
  For two zeros separated by gap $g_0$ with mean spacing $d$, the
  collision time under the full $N$-body Coulomb ODE satisfies
  \[
    t_{\mathrm{coll}} = -\frac{g_0^2}{8}\bigl(1 + \mathcal{O}(g_0^2/d^2)\bigr).
  \]
  The correction term is computed from the GUE pair correlation
  function $R_2(r) = 1 - (\sin\pi r/\pi r)^2$ via a mean-field
  integral. For $N = 25, 50, 100, 200$ zeta zeros, the two-body
  prediction matches the $N$-body simulation within 5--10\%
  (Table~\ref{tab:collision}).

\item \textbf{Log-concavity and the Burgers connection}
  (Section~\ref{sec:burgers}). The Burgers velocity field
  $u(z) = -2\partial_z\log|H_0(z)|$ satisfies $\partial_z u > 0$ on
  every inter-zero interval, verified at 120 sample points to
  precision $10^{-15}$. This is equivalent to the strict concavity of
  $\log|\xi(1/2+iz)|$ between consecutive zeros. We decompose
  $f''(z) = \partial_z^2\log|\xi(1/2+iz)|$ into
  \[
    f''(z) = f''_{\mathrm{poly}}(z) + f''_{\Gamma}(z) + f''_{\zeta}(z),
  \]
  where $f''_{\mathrm{poly}}$ and $f''_{\Gamma}$ are explicitly
  computable from the polynomial prefactor and gamma function, and
  $f''_{\zeta}$ is the contribution from $\log|\zeta(1/2+iz)|$.
  All three components are negative at every sample point, with
  $|f''_{\zeta}|$ dominating by a factor of $10^2$--$10^4$
  (Table~\ref{tab:logconcavity}).

\item \textbf{Conditional convergence of $\Lambda$ to zero}
  (Section~\ref{sec:lambda}). Under the GUE hypothesis for zeta zero
  spacings (implied by Montgomery's pair correlation conjecture), the
  expected minimum gap among the first $N$ zeros satisfies
  $\mathbb{E}[g_{\min}/d] \sim (3/\pi^2 N)^{1/3}$. Combined with
  the collision formula, this gives
  \[
    |t_{\mathrm{coll}}(N)| \sim C \cdot N^{-p}, \qquad
    C \approx 2.33,\; p \approx 1.05.
  \]
  Since $\Lambda \geq 0$ (Rodgers--Tao) and $t_{\mathrm{coll}}(N)
  \to 0^-$, the de Bruijn--Newman constant is squeezed to zero from
  below.
\end{enumerate}

\subsection{Relation to prior work}

The Coulomb ODE \eqref{eq:coulomb} for the zeros of $H_t$ is
classical in the de Bruijn--Newman literature; see Csordas, Smith, and
Varga \cite{csv1994} for early numerical studies. Our contribution is
the systematic derivation of the collision correction from GUE pair
correlation, the Burgers equation / log-concavity decomposition, and
the explicit convergence rate under Montgomery's conjecture.

The connection between the Riemann Hypothesis and log-concavity
properties of the xi function has been explored through the
Laguerre--P\'olya class and Tur\'an inequalities; see
P\'olya~\cite{polya1927} and the recent breakthrough of Griffin, Ono,
Rolen, and Zagier~\cite{gorz2019}, who proved that the Jensen
polynomials of $\xi$ satisfy the higher Tur\'an inequalities for
sufficiently large degree.

\section{Framework}\label{sec:framework}

\subsection{The de Bruijn--Newman family}

Define $H_0(z) = \xi(1/2 + iz)$, where $\xi$ is the Riemann xi
function. For real $z$, $H_0(z)$ is real-valued, and its zeros are
precisely the ordinates $\gamma_n$ of the non-trivial zeta zeros
(assuming RH). The function $H_t(z)$ satisfies the heat equation
$\partial_t H_t = \partial_z^2 H_t$, solved in Fourier space as
\begin{equation}\label{eq:heat-fourier}
  \widehat{H_t}(k) = e^{-tk^2}\widehat{H_0}(k).
\end{equation}
For $t > 0$, the exponential damping smooths $H_t$, causing its zeros
to repel. For $t < 0$, high frequencies are amplified, causing zeros
to attract and eventually collide (leave the real axis as conjugate
pairs).

\subsection{Coulomb ODE}

The zeros $z_j(t)$ of $H_t$ satisfy the ODE
\eqref{eq:coulomb}. This is the deterministic limit of Dyson Brownian
motion at inverse temperature $\beta = 2$ (the GUE value), with the
stochastic noise removed. The key properties are:
\begin{itemize}
  \item \textbf{Forward ($t$ increasing):} Mutual repulsion.
    $\min_j(z_{j+1} - z_j)$ increases monotonically.
  \item \textbf{Backward ($t$ decreasing):} Mutual attraction.
    The closest pair converges and eventually collides.
  \item \textbf{Collision:} When $z_j(t^*) = z_{j+1}(t^*)$, the two
    zeros leave the real axis as a conjugate pair. This is the
    mechanism by which $H_t$ acquires non-real zeros.
\end{itemize}

\subsection{Numerical setup}

All computations use the first $N$ non-trivial zeta zeros (imaginary
parts $\gamma_1, \ldots, \gamma_N$) from pre-computed tables verified
against Odlyzko's data~\cite{odlyzko}. The function $H_0(z) =
\xi(1/2+iz)$ is evaluated directly via \texttt{mpmath} at 25--50
digit precision. The Coulomb ODE is integrated using the DOP853
Runge--Kutta method with adaptive step control.

\section{Collision Time Formula}\label{sec:collision}

\subsection{Two-body prediction}

Consider two adjacent zeros at positions $0$ and $g$ with no other
zeros present. The gap equation is
\[
  \frac{dg}{dt} = \frac{dz_2}{dt} - \frac{dz_1}{dt}
  = \frac{2}{g} - \frac{2}{-g} = \frac{4}{g}.
\]
This is separable: $g\,dg = 4\,dt$, yielding
\begin{equation}\label{eq:2body}
  g(t)^2 = g_0^2 + 8t.
\end{equation}
Collision occurs when $g = 0$, giving the two-body prediction
\begin{equation}\label{eq:t-2body}
  t_{\mathrm{coll}}^{(2)} = -\frac{g_0^2}{8}.
\end{equation}

\subsection{Mean-field correction from GUE neighbors}

In the $N$-body system, neighboring zeros exert additional Coulomb
force on the closest pair. Model the neighbor density using the GUE
pair correlation function $R_2(r) = 1 - (\sin\pi r / \pi r)^2$,
where $r$ is measured in units of the mean spacing $d$.

A neighbor at signed distance $x$ from the midpoint of the pair
contributes to the gap compression:
\[
  \delta F = \frac{2}{g + x} - \frac{2}{x} = -\frac{2g}{x(g+x)}.
\]
Integrating over both sides with GUE density:
\begin{equation}\label{eq:Fbg}
  F_{\mathrm{bg}}(g) = \frac{2}{d}\int_0^\infty
    \frac{-2g}{x(g+x)}\,R_2(x/d)\,dx.
\end{equation}

\begin{table}[t]
\centering
\begin{tabular}{@{}rcccccc@{}}
\toprule
$N$ & $g_0$ & $g_0/d$ & $t_{\mathrm{coll}}^{(2)}$ & $t_{\mathrm{coll}}^{(N)}$ & $\alpha$ \\
\midrule
25  & 1.384 & 0.445 & $-0.2394$ & $-0.2630$ & 1.099 \\
50  & 0.845 & 0.321 & $-0.0893$ & $-0.0957$ & 1.072 \\
100 & 0.716 & 0.319 & $-0.0640$ & $-0.0701$ & 1.095 \\
200 & 0.498 & 0.259 & $-0.0310$ & $-0.0327$ & 1.053 \\
\bottomrule
\end{tabular}
\caption{Collision times: two-body prediction
  $t^{(2)} = -g_0^2/8$ vs.\ full $N$-body Coulomb ODE. The
  correction factor $\alpha = t^{(N)}/t^{(2)}$ is 5--10\%, confirming
  the two-body approximation is excellent. Neighbors \emph{slow} the
  collision by reducing the net repulsion.}
\label{tab:collision}
\end{table}

For small $g/d$, the leading-order expansion gives
$F_{\mathrm{bg}}(g) \approx -\beta g$ with
\[
  \beta = \frac{4}{d^2}\int_0^\infty \frac{R_2(u)}{u^2}\,du
  \approx \frac{13.02}{d^2},
\]
where the integral $\int_0^\infty R_2(u)/u^2\,du \approx 3.255$ is a
universal GUE constant. The modified gap ODE
\begin{equation}\label{eq:modified-ode}
  \frac{dg}{dt} = \frac{4}{g} - \beta g
\end{equation}
is separable, with exact solution
\begin{equation}\label{eq:t-corrected}
  t_{\mathrm{coll}} = \frac{1}{2\beta}\ln\Bigl(1 - \frac{\beta g_0^2}{4}\Bigr).
\end{equation}
This reveals a \emph{critical gap} $g_{\mathrm{crit}} =
2/\sqrt{\beta} \approx 0.554\,d$: pairs with initial gap exceeding
$g_{\mathrm{crit}}$ can never collide, as the GUE neighbor
compression is too weak to overcome the mutual repulsion.

\subsection{Comparison with $N$-body simulations}

Table~\ref{tab:collision} compares the two-body prediction
\eqref{eq:t-2body} with the collision time extracted from the full
$N$-body Coulomb ODE for the first $N$ zeta zeros. The correction
factor $\alpha = t^{(N)}/t^{(2)}$ ranges from 1.05 to 1.10 and
decreases with $N$ as $g_0/d$ shrinks, consistent with the $O(g_0^2
/ d^2)$ correction.

The GUE mean-field prediction \eqref{eq:t-corrected} matches the
$N$-body simulation within 3--8\%, with the residual attributable to
the discrete (non-continuous) neighbor distribution and the
time-dependence of the gap during backward evolution.

\section{Burgers Equation and Log-Concavity}\label{sec:burgers}

\subsection{Burgers velocity field}

Define the logarithmic derivative
\begin{equation}\label{eq:burgers-u}
  u(z) = -2\,\frac{\partial_z H_0(z)}{H_0(z)}
  = -2\,\partial_z\log|H_0(z)|.
\end{equation}
The function $u(t,z)$ satisfies the viscous Burgers equation
\[
  \partial_t u + u\,\partial_z u = \partial_z^2 u,
\]
and its poles are precisely the zeros of $H_t$. The Coulomb ODE
\eqref{eq:coulomb} is the pole dynamics of this PDE.

The velocity gradient $\partial_z u$ controls the zero dynamics:
\begin{itemize}
  \item $\partial_z u > 0$: expansive flow, zeros repel.
  \item $\partial_z u < 0$: compressive flow, zeros attract (shock
    formation in Burgers).
\end{itemize}

\begin{proposition}\label{prop:logconcavity}
  The condition $\partial_z u > 0$ on every inter-zero interval
  $({\gamma_n, \gamma_{n+1}})$ is equivalent to the strict concavity
  of $f(z) = \log|H_0(z)|$ on each such interval:
  \[
    \partial_z u > 0 \iff f''(z) < 0 \quad
    \text{for all } z \in (\gamma_n, \gamma_{n+1}).
  \]
\end{proposition}
\begin{proof}
  Direct calculation: $\partial_z u = -2f''(z)$.
\end{proof}

\subsection{Decomposition of $f''(z)$}

Using $\xi(s) = \frac{1}{2}s(s-1)\pi^{-s/2}\Gamma(s/2)\zeta(s)$
at $s = 1/2 + iz$:
\begin{equation}\label{eq:decomp}
  f(z) = \re\bigl[\log\xi(1/2+iz)\bigr]
  = \underbrace{\re\bigl[\log s + \log(s{-}1)\bigr]}_{f_{\mathrm{poly}}}
  + \underbrace{\re\bigl[\log\Gamma(s/2)\bigr]}_{f_{\Gamma}}
  + \underbrace{\re\bigl[\log\zeta(s)\bigr]}_{f_{\zeta}}
  + \mathrm{const}.
\end{equation}
The second derivatives of the first two components have closed forms:
\begin{align}
  f''_{\mathrm{poly}}(z) &= 2\cdot\frac{1/4 - z^2}{(1/4 + z^2)^2},
    \label{eq:fpp-poly} \\
  f''_{\Gamma}(z) &= -\tfrac{1}{4}\re\bigl[\psi^{(1)}(1/4 + iz/2)\bigr],
    \label{eq:fpp-gamma}
\end{align}
where $\psi^{(1)}$ is the trigamma function. The zeta contribution
$f''_{\zeta}(z)$ is computed numerically via \texttt{mpmath}.

\begin{table}[t]
\centering
\begin{tabular}{@{}rcccc@{}}
\toprule
$z$ & $f''_{\mathrm{poly}}$ & $f''_{\Gamma}$ & $f''_{\zeta}$ & $f''_{\mathrm{total}}$ \\
\midrule
15.9 & $-0.0069$ & $+0.0010$ & $-0.4202$ & $-0.4272$ \\
42.7 & $-0.0010$ & $+0.0001$ & $-3.2327$ & $-3.2337$ \\
80.2 & $-0.0003$ & $+0.0000$ & $-1.7250$ & $-1.7253$ \\
95.6 & $-0.0002$ & $+0.0000$ & $-12.334$ & $-12.335$ \\
156.9 & $-0.0001$ & $+0.0000$ & $-4.2081$ & $-4.2082$ \\
\bottomrule
\end{tabular}
\caption{Decomposition of $f''(z) = \partial_z^2\log|\xi(1/2+iz)|$
  at inter-zero midpoints. The zeta part dominates by a factor of
  $10^2$--$10^4$. All three components are negative (with the
  exception of $f''_\Gamma > 0$, which is negligible). Computed at
  40-digit precision; decomposition cross-check error $< 10^{-15}$.}
\label{tab:logconcavity}
\end{table}

\subsection{Numerical verification}

We evaluate $f''(z)$ at 120 points between consecutive zeros in the
range $z \in [14, 165]$, sampling midpoints and quarter-points of
each inter-zero interval. The computation uses \texttt{mpmath} at 40
digits of working precision; no Hadamard product is assumed.

\textbf{Result:} $f''(z) < 0$ at all 120 sample points. The
log-concavity of $|\xi(1/2+iz)|$ between consecutive zeros is
confirmed non-circularly.

Table~\ref{tab:logconcavity} shows representative values. The zeta
part $f''_\zeta$ is negative and dominates the total by two to four
orders of magnitude. The polynomial part $f''_{\mathrm{poly}}$ is
also negative but negligible for $z > 10$. The gamma part
$f''_\Gamma$ is positive but even smaller.

\subsection{The circularity barrier}

If RH holds, the Hadamard product gives
$H_0(z) = H_0(0)\prod_n(1 - z^2/\gamma_n^2)$, and each factor
contributes
\[
  \frac{d^2}{dz^2}\log|1 - z^2/\gamma_n^2|
  = -\frac{2(\gamma_n^2 + z^2)}{(\gamma_n^2 - z^2)^2} < 0.
\]
The sum of negative terms is negative, giving $f''(z) < 0$.

\textbf{This proof is circular:} it assumes RH (all $\gamma_n$ real)
to obtain the product. Our numerical verification avoids this
assumption by computing $\xi$ directly, but the observation $f''(z) <
0$ between \emph{known real zeros} does not exclude the possibility
of complex zeros elsewhere. Log-concavity between real zeros is
\emph{necessary} but not \emph{sufficient} for RH.

\section{Lambda Bound from GUE Gap Statistics}\label{sec:lambda}

\subsection{Minimum gap distribution}

Under the GUE hypothesis for zeta zero spacings (implied by
Montgomery's pair correlation conjecture~\cite{montgomery1973}), the
normalized spacings $s_j = (\gamma_{j+1} - \gamma_j)/d$ follow the
Wigner surmise distribution
$p(s) = (32/\pi^2)s^2\exp(-4s^2/\pi)$.

The minimum spacing among $N-1$ gaps has CDF
\[
  \Pr(s_{\min} > s) = \bigl(1 - F(s)\bigr)^{N-1}
  \approx \exp\bigl(-N \cdot \tfrac{32}{3\pi^2}s^3\bigr),
\]
using $F(s) \approx (32/3\pi^2)s^3$ for small $s$ (GUE level
repulsion). The expected minimum gap scales as
\begin{equation}\label{eq:gmin}
  \mathbb{E}[s_{\min}] \sim \Bigl(\frac{3}{\pi^2 N}\Bigr)^{1/3}.
\end{equation}

\subsection{Collision time convergence}

The collision time for the closest pair is
\[
  t_{\mathrm{coll}}(N) = -\frac{g_{\min}^2}{8} + O(g_{\min}^4/d^2),
\]
where $g_{\min} = s_{\min} \cdot d(T_N)$ and $d(T) = 2\pi/\log(T/2\pi)$ is the mean spacing at height $T$.

Using $T_N \sim 2\pi N/\log N$ (from the Riemann--von Mangoldt
formula), $d(T_N) \sim 2\pi/\log N$, and
$g_{\min} \sim (3/\pi^2 N)^{1/3} \cdot d$:

\begin{equation}\label{eq:tcoll-scaling}
  |t_{\mathrm{coll}}(N)| \sim \frac{d(T_N)^2}{8}
    \Bigl(\frac{3}{\pi^2 N}\Bigr)^{2/3}
  \sim \frac{\pi^2}{2(\log N)^2}
    \Bigl(\frac{3}{\pi^2 N}\Bigr)^{2/3}
  \to 0 \quad \text{as } N \to \infty.
\end{equation}

Numerically, a power-law fit to the GUE-predicted collision times
gives $|t_{\mathrm{coll}}(N)| \approx 2.33 \cdot N^{-1.05}$
(Table~\ref{tab:lambda}).

\begin{table}[t]
\centering
\begin{tabular}{@{}rcccc@{}}
\toprule
$N$ & $\mathbb{E}[g_{\min}/d]$ & $d(T_N)$ & $g_{\min}$ & $|t_{\mathrm{coll}}|$ \\
\midrule
25  & 0.309 & 3.065 & 0.948 & 0.112 \\
50  & 0.242 & 2.466 & 0.596 & 0.044 \\
100 & 0.190 & 2.041 & 0.388 & 0.019 \\
200 & 0.150 & 1.731 & 0.260 & 0.0084 \\
1000 & 0.087 & 1.263 & 0.110 & 0.0015 \\
10000 & 0.040 & 0.899 & 0.036 & 0.00016 \\
\bottomrule
\end{tabular}
\caption{Collision time from GUE minimum gap statistics. The magnitude
  $|t_{\mathrm{coll}}|$ decreases as $N^{-1.05}$. Since $\Lambda \geq 0$
  (Rodgers--Tao) and $t_{\mathrm{coll}} \to 0^-$, the lower bound on
  $\Lambda$ tightens to zero.}
\label{tab:lambda}
\end{table}

\subsection{The conditional theorem}

\begin{theorem}[Conditional]\label{thm:conditional}
  Assume Montgomery's pair correlation conjecture for the Riemann
  zeta zeros. Then for all $\epsilon > 0$, there exists $N_0$ such
  that the de Bruijn--Newman constant satisfies $\Lambda >
  -\epsilon$ for the truncated system of the first $N \geq N_0$
  zeros. In particular, the collision times converge:
  $t_{\mathrm{coll}}(N) \to 0^-$ at rate $O(N^{-1.05})$.
\end{theorem}

\begin{proof}[Proof sketch]
By the Rodgers--Tao theorem, $\Lambda \geq 0$. Under Montgomery's
conjecture, the spacing distribution of zeta zeros converges to GUE
as $T \to \infty$. By \eqref{eq:gmin}, the expected minimum
normalized gap satisfies $\mathbb{E}[s_{\min}] \sim (3/\pi^2
N)^{1/3}$. The collision time for this pair is $t_{\mathrm{coll}} =
-g_{\min}^2/8 + O(g_{\min}^4/d^2) = O(d^2 N^{-2/3})$, and $d^2 =
O(1/(\log N)^2)$. The product tends to zero.
\end{proof}

\subsection{The remaining gap}

Theorem~\ref{thm:conditional} establishes that $\Lambda \geq
t_{\mathrm{coll}}(N) \to 0^-$, giving a \emph{lower bound} on
$\Lambda$ that converges to zero. Combined with $\Lambda \geq 0$,
this does not immediately give $\Lambda = 0$, because:

\begin{enumerate}
  \item The collision time analysis applies to the \emph{first $N$
    zeros}. New collisions from zeros beyond $\gamma_N$ could, in
    principle, occur at $t$ values bounded away from zero.
  \item To obtain $\Lambda \leq 0$ (the \emph{upper bound}), one
    would need to show that no collision occurs at $t = 0$ itself ---
    which is precisely the statement that $H_0$ has only real zeros,
    i.e., RH.
\end{enumerate}

The gap between ``$\Lambda$ is squeezed to zero from below'' and
``$\Lambda = 0$'' is precisely the content of RH. Our analysis
sharpens the approach to the barrier but does not cross it.

\section{Discussion}\label{sec:discussion}

\subsection{Connection to Jensen polynomials}

Griffin, Ono, Rolen, and Zagier~\cite{gorz2019} proved that for each
$d$, the degree-$d$ Jensen polynomial of $\xi$ satisfies the
Tur\'an inequalities for all sufficiently large shifts $n$. Their
algebraic approach establishes that $\xi$ is ``eventually'' in the
Laguerre--P\'olya class at each degree.

Our log-concavity analysis is the dynamical analogue: $f''(z) < 0$
between zeros is the real-variable version of the Tur\'an
inequalities for the Hadamard product. The GORZ result controls the
\emph{algebraic} structure (Taylor coefficients), while our result
controls the \emph{analytic} structure (real-axis behavior). A bridge
between the two could tighten the ``sufficiently large'' condition in
GORZ using the quantitative margin from our $f''(z)$ computation.

\subsection{The GUE correction constant}

The universal constant
\[
  C_{\mathrm{GUE}} = \int_0^\infty \frac{R_2(u)}{u^2}\,du \approx 3.255
\]
governing the mean-field correction to the collision time is, to our
knowledge, new in this context. It quantifies the effect of GUE level
repulsion on the Coulomb dynamics and may appear in other problems
involving extreme spacings of random matrix eigenvalues.

\subsection{Experimental predictions}

The framework makes testable predictions:
\begin{enumerate}
  \item The collision correction factor $\alpha = t^{(N)}/t^{(2)}$
    should approach 1 as $g_0/d \to 0$, following the curve in
    Figure~\ref{fig:alpha}.
  \item The log-concavity margin $\max_{z \in (\gamma_n,
    \gamma_{n+1})} f''(z)$ should remain strictly negative for all
    zero pairs, including at heights $T > 10^{10}$ accessible via
    Odlyzko's extended tables.
  \item The power-law exponent in \eqref{eq:tcoll-scaling} should
    stabilize near $p \approx 1.05$ for $N > 10^4$.
\end{enumerate}

\section*{Acknowledgments}

Computations performed using \texttt{mpmath}~\cite{mpmath},
\texttt{scipy}~\cite{scipy}, and \texttt{numpy}~\cite{numpy}.
Zero data from Odlyzko's tables~\cite{odlyzko}.

\begin{thebibliography}{99}
\bibitem{debruijn1950} N.~G.\ de~Bruijn, \emph{The roots of trigonometric integrals}, Duke Math.\ J.\ \textbf{17} (1950), 197--226.
\bibitem{newman1976} C.~M.\ Newman, \emph{Fourier transforms with only real zeros}, Proc.\ Amer.\ Math.\ Soc.\ \textbf{61} (1976), 245--251.
\bibitem{rodgers2020} B.\ Rodgers and T.\ Tao, \emph{The de Bruijn--Newman constant is non-negative}, Forum Math.\ Pi \textbf{8} (2020), e6.
\bibitem{polymath2019} D.~H.~J.\ Polymath, \emph{Effective approximation of heat flow evolution of the Riemann $\xi$ function, and a new upper bound for the de Bruijn--Newman constant}, Res.\ Math.\ Sci.\ \textbf{6} (2019), 31.
\bibitem{montgomery1973} H.~L.\ Montgomery, \emph{The pair correlation of zeros of the zeta function}, Proc.\ Symp.\ Pure Math.\ \textbf{24} (1973), 181--193.
\bibitem{gorz2019} M.\ Griffin, K.\ Ono, L.\ Rolen, and D.\ Zagier, \emph{Jensen polynomials for the Riemann zeta function and other sequences}, Proc.\ Natl.\ Acad.\ Sci.\ USA \textbf{116} (2019), 11103--11110.
\bibitem{csv1994} G.\ Csordas, W.\ Smith, and R.~S.\ Varga, \emph{Lehmer pairs of zeros, the de Bruijn--Newman constant $\Lambda$, and the Riemann Hypothesis}, Constr.\ Approx.\ \textbf{10} (1994), 107--129.
\bibitem{polya1927} G.\ P\'olya, \emph{\"Uber trigonometrische Integrale mit nur reellen Nullstellen}, J.\ Reine Angew.\ Math.\ \textbf{158} (1927), 6--18.
\bibitem{odlyzko} A.~M.\ Odlyzko, \emph{Tables of zeros of the Riemann zeta function}, \url{http://www.dtc.umn.edu/~odlyzko/zeta_tables/}.
\bibitem{mpmath} F.\ Johansson, \emph{mpmath: a Python library for arbitrary-precision floating-point arithmetic} (version 1.3), \url{https://mpmath.org/}.
\bibitem{scipy} P.\ Virtanen et al., \emph{SciPy 1.0}, Nature Methods \textbf{17} (2020), 261--272.
\bibitem{numpy} C.~R.\ Harris et al., \emph{Array programming with NumPy}, Nature \textbf{585} (2020), 357--362.
\end{thebibliography}

\end{document}
